\silentchapter{Foundation}{foundation}
\begin{widebody}
{\small

The framework is a heuristic. It does not say what the form will be; it says what one must look for. Such is all operational science: Pavlovian learning is a heuristic that can be operationalised at every level from gene regulatory networks to civilisations, and its value lies in making the same question tractable across scales. The free-energy principle\textsuperscript{19,\,83} does the same for agency. The principle is, as of 2026, empirically demonstrated across at least fourteen substrate levels. An excerpt shows the range.

\bigskip
\begin{center}
\footnotesize
\renewcommand{\arraystretch}{1.2}
\begin{tabular}{@{}p{38mm}p{50mm}l@{}}
\toprule
\textbf{Scale} & \textbf{Mechanism} & \textbf{Source} \\
\midrule
Gene regulatory network & Six Pavlovian learning forms & \textsuperscript{36} \\
Single cell & Habituation, associative conditioning & \textsuperscript{98} \\
Slime mould & Habituation, memory transfer & \textsuperscript{99} \\
Tissue & Bioelectric morphogenetic memory & \textsuperscript{85} \\
Immune system & Trained immunity via epigenetics & \textsuperscript{100} \\
Brain & Synaptic plasticity, predictive processing & \textsuperscript{101} \\
Individual & Six Pavlovian forms, symbolic learning & \textsuperscript{102} \\
Dyad & Bayesian synchrony & \textsuperscript{103} \\
Institution & Exploration and exploitation & \textsuperscript{104} \\
Civilisation & Paradigm shift as collective relearning & \textsuperscript{105} \\
\bottomrule
\end{tabular}
\end{center}
\bigskip

That the same formalism predicts phenomena across a scale from nanometer to continent is an unreasonable effectiveness in Wigner's\textsuperscript{106} sense, \textit{the unreasonable effectiveness of mathematics}. It justifies pragmatic use even before the deepest interpretations have been settled.

The effectiveness is visible outside mathematics too. Machine learning has made style classification a solved problem; what to a design historian is a life's work, a calibrated neural network can perform in milliseconds with higher agreement than the expert panel. The parallel exists in medicine. Skin cancer classification is now solved in the sense that machines systematically surpass experienced dermatologists in sensitivity and specificity; not to use them would be irresponsible. The finding is not that the clinician is replaced. The clinician works with the machine and spends the freed time on what the machine does not do: clinical conversations, complicated cases, boundary surfaces. The machine handles the patterned; the human handles the pattern-breaking.

The design theorist's situation is homologous. So long as most quantitative design research is style classification, which is a category mistake (cf. 2.212) and a solved problem, the field ties its research time to what the machine can already do. The framework does not replace the design theorist; it frees the time for what no machine can do: operationalisation of care, identification of unrepresented axes, articulation of a research programme that demands domain insight.

An observation on what kind of progress this is. The history of science over the past century can be read as a systematic dissolution of binary categories. Every time a hard dichotomy was reckoned as a heuristic rather than an ontology, it opened a research space that had been closed.

\begin{center}
\footnotesize
\setlength{\tabcolsep}{5pt}
\renewcommand{\arraystretch}{1.4}
\arrayrulecolor{gray!50}
\begin{tabular}{@{}p{42mm}p{56mm}@{}}
\toprule
\textbf{Binary category} & \textbf{The heuristic that dissolved it} \\
\midrule
Wave and particle & Duality, quantum mechanics \\
Matter and energy & $E=mc^2$, unified physics \\
Gene and environment & Epigenetics, developmental biology \\
Living and non-living & Autopoiesis, FEP as continuum \\
Brain and body & Embodied cognition, interstitium \\
Species and species & Hybridisation, horizontal gene transfer \\
Nature and culture & Niche construction, multiscale evolution \\
Human and animal & Competency scale, cognitive continuity \\
Consciousness and non-consciousness & Graded agency, multiscale competency \\
Object and relation & Interstitium, systems biology, network science \\
\bottomrule
\end{tabular}
\arrayrulecolor{black}
\end{center}

Design theory stands in this lineage. It stands there with a historical debt: design thinking has for decades broken from the academy and found a home elsewhere, because design theory lacked the foundation to hold it. A discipline cannot remain rootless for long. The binary distinctions that still hold the design field in place are form and function, nature and artefact, intention and emergence, biological and cultural, human and technological, form-maker and user, designer and form-maker. Each is a concrete research programme: find the heuristic that turns the distinction into a gradient, and measure what research space opens. To treat these as gradients is not merely to borrow a move from the natural sciences; it is to follow the only movement that has produced lasting progress.

The act of design is formally isomorphic with active inference: an agent with a generative model of a form to be, an observation of the form that is, and an adjustment mechanism between them. The formalism is substrate-independent. The propositions of the tractatus therefore offer a common language in which the act of design can be integrated directly with machine learning (via the generative model in 8.1), developmental biology (via substrate-independent agents), ecology (via the multiscale hierarchy), and climate science (via the research ethics that follows from section 7). The tools we think with are extended.

When the light cone and care are substrate-independent, the act of design becomes applicable to domains it has been foreign to. To shape a nature reserve is to navigate a morphospace in which microbes, fungi, birds, moose, and climate zones are all form-makers; the designer is an agent among them, not above. To shape an AI agent is to design a new form-maker, not a tool; scaling, light cone, and care are design parameters. Each such extension opens a channel to a field that today works without the designer, from biosciences to infrastructure planning. The economic argument follows by itself: the field as it is trains candidates who build targeting systems and dependency interfaces, and these are net costs to society. The field as it could be trains candidates who work with microbiologists, ecologists, climate scientists, AI researchers; net benefits, measurable in concrete dividends in the fields they work with. The ratio between the two sums is several orders of magnitude.

This gives the field a research programme in four points. \textit{First}, to define the core concepts operationally, with reference to a framework that is already validated. Form, form-maker, and designer are no longer freely defined; they are positions in an adaptation landscape, agents that navigate it, and species of agent with specific target states. \textit{Second}, to operationalise care as the set of axes an agent actively adjusts for, and to develop quantitative measures; the K metric\textsuperscript{86} provides one. \textit{Third}, to take up the enormous empirical archives that museums, libraries, ethnographers, and statisticians have gathered over generations. These archives are not scarce; they are underutilised, often interpreted with category mistakes where style is treated as cause, function as primary pressure, and the user as the only agent. Reanalysis of existing design research data under this lens is a concrete research programme. \textit{Fourth}, to bind professional certification to documented breadth of care, as medicine binds certification to documented effect.

The material is there. Formlære's contribution is to bind it together.

}
\end{widebody}
